% GENERATED FILE. Edit canonical lessons, then run npm run generate:books.
% canonical-source-hash: fnv1a64:9b75c3201b394318
% canonical-lessons: TA-C133-tol, TA-C133-marpu, TA-C133-mulankal, TA-C133-muti, TA-C133-tol2

\chapter{Shoulder, Chest, Knee, Hair, Skin}
\label{ch:ta-shoulder-chest-knee-hair-skin}
\hlchaptermodality{133}

\begin{chapteropening}
\textbf{By the end of this chapter:} \emph{I can say a shoulder, the chest, a knee, hair and the skin.}

Complete the last lesson of chapter 133: i can say a shoulder, the chest, a knee, hair and the skin.
\end{chapteropening}

\section[tōḷ]{\ta{தோள்} (tōḷ) — a shoulder}

\label{lesson:TA-C133-tol}



\begin{quote}
Before the new one: say the Tamil for the tongue, then the Tamil for a finger.
\end{quote}

\subsection*{You'll want to know: \ta{தோள்}}
\textbf{\ta{தோள்}} — \emph{tōḷ} — \textquotedblleft{}a shoulder\textquotedblright{}.

\begin{grammarlens}[title={What you've built}]
One more word for this chapter.
\end{grammarlens}

\subsection*{Your turn}
\begin{itemize}
\raggedright
  \item \emph{Say it:} \emph{tōḷ}
  \item \emph{Say it:} \emph{tōḷ}, once more
  \item \emph{Say it:} say \emph{tōḷ}
  \item \emph{Recall:} say the Tamil for the tongue, then the Tamil for a finger, then say \emph{tōḷ} again
\end{itemize}

\begin{tcolorbox}[breakable,colback=teal!4,colframe=teal!35!black,title={Before you move on}]
What does \ta{தோள்} mean? (\textquotedblleft{}A shoulder\textquotedblright{}.) Say it once more.
\end{tcolorbox}
\section[mārpu]{\ta{மார்பு} (mārpu) — the chest}

\label{lesson:TA-C133-marpu}



\begin{quote}
Before the new one: say the Tamil for a finger, then the Tamil for a shoulder.
\end{quote}

\subsection*{You'll want to know: \ta{மார்பு}}
\textbf{\ta{மார்பு}} — \emph{mārpu} — \textquotedblleft{}the chest\textquotedblright{}.

\begin{grammarlens}[title={What you've built}]
One more word for this chapter.
\end{grammarlens}

\subsection*{Your turn}
\begin{itemize}
\raggedright
  \item \emph{Say it:} \emph{mārpu}
  \item \emph{Say it:} \emph{mārpu}, once more
  \item \emph{Say it:} say \emph{mārpu}
  \item \emph{Recall:} say the Tamil for a finger, then the Tamil for a shoulder, then say \emph{mārpu} again
\end{itemize}

\begin{tcolorbox}[breakable,colback=teal!4,colframe=teal!35!black,title={Before you move on}]
What does \ta{மார்பு} mean? (\textquotedblleft{}The chest\textquotedblright{}.) Say it once more.
\end{tcolorbox}
\section[muḻaṅkāl]{\ta{முழங்கால்} (muḻaṅkāl) — a knee}

\label{lesson:TA-C133-mulankal}



\begin{quote}
Before the new one: say the Tamil for a shoulder, then the Tamil for the chest.
\end{quote}

\subsection*{You'll want to know: \ta{முழங்கால்}}
\textbf{\ta{முழங்கால்}} — \emph{muḻaṅkāl} — \textquotedblleft{}a knee\textquotedblright{}.

\begin{grammarlens}[title={What you've built}]
One more word for this chapter.
\end{grammarlens}

\subsection*{Your turn}
\begin{itemize}
\raggedright
  \item \emph{Say it:} \emph{muḻaṅkāl}
  \item \emph{Say it:} \emph{muḻaṅkāl}, once more
  \item \emph{Say it:} say \emph{muḻaṅkāl}
  \item \emph{Recall:} say the Tamil for a shoulder, then the Tamil for the chest, then say \emph{muḻaṅkāl} again
\end{itemize}

\begin{tcolorbox}[breakable,colback=teal!4,colframe=teal!35!black,title={Before you move on}]
What does \ta{முழங்கால்} mean? (\textquotedblleft{}A knee\textquotedblright{}.) Say it once more.
\end{tcolorbox}
\section[muṭi]{\ta{முடி} (muṭi) — hair}

\label{lesson:TA-C133-muti}



\begin{quote}
Before the new one: say the Tamil for the chest, then the Tamil for a knee.
\end{quote}

\subsection*{You'll want to know: \ta{முடி}}
\textbf{\ta{முடி}} — \emph{muṭi} — \textquotedblleft{}hair\textquotedblright{}.

\begin{grammarlens}[title={What you've built}]
One more word for this chapter.
\end{grammarlens}

\subsection*{Your turn}
\begin{itemize}
\raggedright
  \item \emph{Say it:} \emph{muṭi}
  \item \emph{Say it:} \emph{muṭi}, once more
  \item \emph{Say it:} say \emph{muṭi}
  \item \emph{Recall:} say the Tamil for the chest, then the Tamil for a knee, then say \emph{muṭi} again
\end{itemize}

\begin{tcolorbox}[breakable,colback=teal!4,colframe=teal!35!black,title={Before you move on}]
What does \ta{முடி} mean? (\textquotedblleft{}Hair\textquotedblright{}.) Say it once more.
\end{tcolorbox}
\section[tōl]{\ta{தோல்} (tōl) — the skin}

\label{lesson:TA-C133-tol2}



\begin{quote}
Before the new one: say the Tamil for a knee, then the Tamil for hair.
\end{quote}

\subsection*{You'll want to know: \ta{தோல்}}
\textbf{\ta{தோல்}} — \emph{tōl} — \textquotedblleft{}the skin\textquotedblright{}.

\begin{grammarlens}[title={What you've built}]
One more word for this chapter.
\end{grammarlens}

\subsection*{Your turn}
\begin{itemize}
\raggedright
  \item \emph{Say it:} \emph{tōl}
  \item \emph{Say it:} \emph{tōl}, once more
  \item \emph{Say it:} say \emph{tōl}
  \item \emph{Recall:} say the Tamil for a knee, then the Tamil for hair, then say \emph{tōl} again
\end{itemize}

\begin{tcolorbox}[breakable,colback=teal!4,colframe=teal!35!black,title={Before you move on}]
What does \ta{தோல்} mean? (\textquotedblleft{}The skin\textquotedblright{}.) Say it once more.
\end{tcolorbox}
